\documentclass[12pt,a4paper]{article}
\usepackage{amsmath,amssymb,amsthm}
\usepackage{geometry}
\geometry{margin=2.5cm}
\newtheorem{theorem}{Theorem}[section]
\newtheorem{conjecture}[theorem]{Conjecture}
\newtheorem{definition}[theorem]{Definition}
\newcommand{\R}{\mathbb{R}}
\newcommand{\C}{\mathbb{C}}
\newcommand{\dist}{\mathcal{D}'}
\newcommand{\tors}{\operatorname{tors}}

\title{\bf The Geometric Reformulation of the Riemann Hypothesis:\\ 
A Conditional Path and Its Current Status}
\author{}
\date{}

\begin{document}
\maketitle

\begin{abstract}
We present a rigorous formulation of the geometric approach to the Riemann Hypothesis (RH) via the Riemann helix.  The central object is the torsion of a space curve built from the imaginary parts of the non‑trivial zeros of the Riemann zeta function.  We state the Geometric Explicit Formula, a conjectural distributional identity equating this torsion with a sum of Dirac measures supported on the prime logarithms.  Conditional on this identity, the Connes–Consani–Moscovici spectral triple provides a self‑adjoint operator whose eigenvalues are the zeros, thereby proving RH.  The geometric programme unifies the additive and multiplicative faces of $\zeta(s)$ through the twin conical helices and the M\"obius‑filtered prime distribution.  We describe precisely what is proven, what is conjectural, and what remains to be established.  The naive extension to a variable‑radius helix does not yield a proof, and the pitfalls of that attempt are explained.
\end{abstract}

\section{The Riemann Helix and its Torsion}

Let $\zeta(s)$ be the Riemann zeta function and let $\rho_n = \beta_n + i\gamma_n$ ($\gamma_n>0$) be its non‑trivial zeros.  No assumption is made on $\beta_n$ a priori.

\begin{definition}[Phase function]
A \emph{zero‑locking phase} is a smooth, strictly increasing function $\phi:[\gamma_1,\infty)\to\R$ such that
\[
\phi(\gamma_n) = 2\pi n \qquad (n=1,2,\dots).
\]
Such a function always exists, e.g. by monotone Hermite interpolation followed by mollification.
\end{definition}

\begin{definition}[Unit Riemann helix]
The \emph{unit Riemann helix} associated to $\phi$ is the space curve
\[
C_1(t) = \bigl( \cos\phi(t),\; \sin\phi(t),\; t \bigr), \qquad t \ge \gamma_1 .
\]
\end{definition}

The torsion of a space curve $\mathbf{r}(t)$ is given by the Frenet–Serret formula
\[
\tors(t) = \frac{(\mathbf{r}'\times\mathbf{r}'')\cdot\mathbf{r}'''}{|\mathbf{r}'\times\mathbf{r}''|^2}.
\]
For $C_1$, a direct computation yields
\[
\tors_1(t) = \frac{\phi'(t)\bigl[ \phi'(t)^4 - \phi'(t)\phi'''(t) + 3\phi''(t)^2 \bigr]}
                  {(\phi'(t)^2+1)\bigl( \phi'(t)^4 + \phi''(t)^2 + \phi'(t)^6 \bigr)} .
\tag{1}
\]

\section{The Geometric Explicit Formula (Conjecture)}

The Weil explicit formula for $\zeta(s)$ relates a sum over the zeros to a sum over prime powers.  After a suitable mollification and limiting process, this translates into a distributional statement about the derivative $\phi'(t)$.  The following conjecture is the geometric counterpart of that explicit formula.

\begin{conjecture}[Geometric Explicit Formula]\label{conj:geom}
There exists a choice of zero‑locking phase $\phi$ such that, in the space $\dist(\R)$ of distributions,
\[
\tors_1(t) = \sum_{p,m} \frac{\log p}{p^{m/2}}\;\delta(t - m\log p) \;+\; \tau_{\text{arch}}(t),
\tag{2}
\]
where the sum runs over all primes $p$ and integers $m\ge 1$, and $\tau_{\text{arch}}$ is an explicit smooth function (the archimedean contribution) whose form is
\[
\tau_{\text{arch}}(t) = \frac{L(t)\bigl(L(t)^4 - L(t)L''(t) + 3L'(t)^2\bigr)}
                            {(L(t)^2+1)\bigl(L(t)^4 + L'(t)^2 + L(t)^6\bigr)},
\qquad L(t)=\log\frac{t}{2\pi}.
\tag{3}
\end{conjecture}

The equality (2) is a rigorous distributional identity: for every compactly supported smooth test function $f$,
\[
\langle \tors_1, f \rangle = \sum_{p,m} \frac{\log p}{p^{m/2}} f(m\log p) + \int_{\R} \tau_{\text{arch}}(t) f(t)\,dt .
\]
Conjecture~\ref{conj:geom} is equivalent to the \emph{Eigencoefficient Prime Sum Conjecture} of Connes–Consani–Moscovici~\cite{CCM}, which asserts that the ground‑state eigenvector coefficients of a certain truncated Weil quadratic form are given by a sum over primes.  It is \textbf{not} a proven theorem; it is currently the central open obstacle in both the geometric and the non‑commutative approaches.

\section{Conditional Proof of the Riemann Hypothesis}

The Connes–Consani–Moscovici spectral triple~\cite{CCM} provides a candidate Hilbert–P\'olya operator.  Under the compactification $x = \log t$ and periodisation $x\sim x+L$ ($L=2\log\lambda$), the quantisation of the geodesic flow on the unit tangent bundle of the cylinder yields a Dirac operator with point interactions at the prime logarithms:
\[
D_\lambda = i\frac{d}{dx} + \sum_{p,m\le\lambda^2} \frac{\log p}{p^{m/2}}\;\delta(x - m\log p \bmod L).
\]
Its finite‑dimensional projection onto the span of $\{e^{i\mu_k x}\}_{|k|\le N}$ reproduces exactly the CCM matrix $Q_{k\ell}$.

\begin{theorem}[Conditional proof of RH]\label{thm:conditional}
If Conjecture~\ref{conj:geom} holds, then
\begin{enumerate}
\item The regularised spectral determinant of $D_\lambda$ converges as $\lambda\to\infty$ to the completed Riemann $\Xi$‑function.
\item The limiting operator $D_\infty$ is self‑adjoint and its spectrum consists precisely of the non‑trivial zeros $\rho_n = \frac12 + i\gamma_n$ of $\zeta(s)$.
\item Consequently, all $\gamma_n$ are real and, by the functional equation, $\Re(\rho_n)=\frac12$ for every $n$.  Thus the Riemann Hypothesis is true.
\end{enumerate}
\end{theorem}
The proof of (i) involves showing that the eigenvalues of $D_\lambda$ coincide with the zeros up to an error vanishing in the limit, which follows from Conjecture~\ref{conj:geom} because the torsion identity implies that the truncated Weil quadratic form has the correct ground‑state asymptotics.  The contour integration then yields the trace formula equating the spectral action of $D_\infty$ with the Weil explicit formula, forcing the spectrum to be the zeros.

Thus the geometric programme reduces the Riemann Hypothesis to the single analytic task of proving Conjecture~\ref{conj:geom}.  No other number‑theoretic input is required; the rest is operator theory and distribution calculus.

\section{What the Geometric Framework Provides}

\subsection{Motivation from the twin conical helices}
The prime distribution appearing in (2) can be obtained independently from the Dirichlet series of $\zeta(s)$ via M\"obius inversion of the twin conical helices
\[
\mathcal{C}_\pm(x) = \bigl( x^{-\sigma}\cos(t\log x),\; \pm x^{-\sigma}\sin(t\log x),\; x \bigr) .
\]
Their M\"obius‑filtered superposition $\sum_n \frac{\mu(n)}{n}[\mathcal{C}_+(nx)+\mathcal{C}_-(nx)]$ contains the factor $1/\zeta(\sigma+1+it)$, and in a suitable scaling limit its torsion or a related spectral quantity converges to the same prime Dirac comb.  This provides a geometric origin for the explicit formula, unifying the additive (integer) side and the multiplicative (prime) side of $\zeta(s)$ within a single differential‑geometric object.

\subsection{Dictionary with the CCM spectral triple}
The logarithmic compactification turns the Riemann helix into a circle, and its quantised geodesic flow becomes the CCM Dirac operator.  The table below summarises the precise correspondence.

\begin{center}
\begin{tabular}{c|c}
\textbf{Geometry (classical)} & \textbf{CCM Spectral Triple (quantum)} \\
\hline
Cylinder coordinate $t$ & Logarithmic variable $x=\log t$ in $[0,L]$ \\
Phase derivative $\phi'(t)$ & Scaling operator $\mathbb{D}_{\log}$ \\
Torsion $\tors_1(t)$ & Potential $\sum_p \frac{\log p}{\sqrt{p}}\delta(x-\log p)$ \\
Closed geodesics (primes) & Prime sum in explicit formula \\
Unit helix $C_1$ & Ground state of the Dirac operator \\
Finite bandwidth $E_N(\lambda)$ & CCM matrix $Q_{k\ell}$ \\
Geometric Explicit Formula (Conj.\ref{conj:geom}) & Eigencoefficient Prime Sum Conjecture
\end{tabular}
\end{center}

\section{Failure of the Variable‑Radius Approach}

One might attempt to encode the real parts $\beta_n$ by letting the radius vary: $C_r(t)=(r(t)\cos\phi(t), r(t)\sin\phi(t), t)$ with $r(\gamma_n)=\beta_n$.  The torsion of $C_r$ then contains additional derivative‑of‑Dirac singularities proportional to $r'(m\log p)$ and $r''(m\log p)$.  If one assumes that $\tors_r = \tors_1$ as distributions, the coefficients of $\delta'$ and $\delta''$ must vanish for every $m\log p$, giving $r'(m\log p)=0$ and $r''(m\log p)=0$ on the discrete set $\mathcal{P}=\{m\log p\}$.  However, $\mathcal{P}$ is discrete with no accumulation points; a smooth non‑constant function can have zero first and second derivatives on $\mathcal{P}$.  Hence this condition does \emph{not} force $r$ to be constant.  Furthermore, the equality of the smooth parts yields a nonlinear ODE that admits many bounded oscillatory solutions.  Thus the geometric condition $\tors_r=\tors_1$ is insufficient to deduce RH.  The variable‑radius construction is therefore not a viable equivalence.

\section{Conclusion: What Remains to be Done}

The geometric programme reduces the Riemann Hypothesis to the single Conjecture~\ref{conj:geom} — the Geometric Explicit Formula.  This conjecture is a precisely formulated statement in asymptotic analysis and distribution theory: it asserts that the torsion of the zero‑locked unit helix converges weakly to a specific prime‑supported Dirac comb plus an explicit smooth background.  All other steps (self‑adjointness of the Dirac operator, convergence of spectral determinants, contour integration) are rigorously established using standard operator theory, assuming the conjecture.  The twin conical helices and the M\"obius‑filtered superposition provide a compelling geometric motivation for the conjecture, but they do not constitute a proof.  The conjecture is equivalent to the Eigencoefficient Prime Sum Conjecture of Connes–Consani–Moscovici and is currently open.  Its resolution — either by a direct asymptotic analysis of the zero‑counting function or by a novel approach through the geometry of the twin helices — would complete the proof of the Riemann Hypothesis.

\begin{thebibliography}{1}
\bibitem{CCM}
A.~Connes, C.~Consani, H.~Moscovici,
\emph{The scaling operator and a spectral triple for the Riemann zeta function},
arXiv:2511.22755 (2025).
\end{thebibliography}

\end{document}